\chapter{Fondements pour la licence : complétude, intégration et abstraction}
\section*{Exercice 36.1 -- Supremum et maximum}
[Base] Déterminer la borne supérieure de $A=\{1-1/n:n\ge1\}$. Est-ce un maximum ?

\section*{Exercice 36.2 -- Une hypothèse indispensable}
[Critique] Une fonction continue sur un intervalle est-elle toujours bornée et atteint-elle toujours ses bornes ?

\section*{Exercice 36.3 -- Convergence uniforme}
[Approfondissement] Pour $f_n(x)=x/n$ sur $[0,1]$, montrer la convergence uniforme vers zéro. Que devient cette affirmation sur $\R$ ?

\section*{Exercice 36.4 -- Noyau d’une dérivation}
[Raisonnement] Sur les polynômes de degré au plus 2, écrire la matrice de la dérivation dans la base $(1,X,X^2)$ et déterminer son noyau.

\section*{Exercice 36.5 -- Échange limite–intégrale}
[Approfondissement] Si $f_n,f$ sont continues sur $[a,b]$ et $f_n\to f$ uniformément, montrer $\int_a^bf_n\to\int_a^bf$.

\section*{Exercice 36.6 -- Audit d’une preuve}
[Synthèse courte] Un élève écrit : « $u_{n+1}=\sqrt{2+u_n}$, donc la limite vérifie $L=\sqrt{2+L}$ et vaut 2. » Identifier les étapes manquantes et compléter pour $u_0=0$.
